\documentclass[a4paper,11pt,fleqn]{article}%fleqn pour aligner les equations à gauche
\usepackage{/home/rweye/Nextcloud/Documents/Overleaf/Styles/Cours}

\newcommand{\chnum}{Chapitre 17}
\newcommand{\chtitle}{Circuits électriques}
\newcommand{\dtype}{Cours}


\begin{document}

\maketitle

\section{Description d'un circuit}
\subsection{Deux types de circuits électriques}
Un circuit électrique est une association de dipôles, comprenant un générateur dont le rôle est de délivrer le courant électrique. Un circuit \textbf{en série} comporte une seule \textbf{maille}. Un circuit \textbf{en dérivation} comporte \textbf{au moins deux mailles}.\\
La portion de circuit entre deux \textbf{nœuds} consécutifs est une \textbf{branche}. La branche principale comporte le générateur et les autres branches sont des branches dérivées.

\vspace{3cm}

\subsection{Le courant électrique}
Le \textbf{courant électrique} circule à l'extérieur du générateur \textbf{de la borne positive vers la borne négative} : c'est le \textbf{sens conventionnel} du courant électrique représenté par une flèche.\\
L'\textbf{intensité $I$ du courant} électrique s'exprime en \textbf{ampères (A)}. Elle se mesure à l'aide d'un \textbf{ampèremètre} branché en série.\\
L'\textbf{intensité} du coutant électrique qui traverse des dipôles associés en série est la même.

\subsection{La tension}
La \textbf{tension $U$} aux bornes d'un dipôle est représentée par une flèche. La tension s'exprime en \textbf{volts (V)}. Elle se mesure avec un \textbf{voltmètre} branché en dérivation aux bornes du dipôle.\\
La \textbf{tension} aux bornes de dipôles associés en \textbf{dérivation} est \textbf{la même}.

\vspace{3cm}

\section{Lois des circuits électriques}
\subsection{La loi des nœuds}
Dans un circuit \textbf{en dérivation}, la somme \textbf{des intensités} des courants électriques qui \textbf{arrivent à un nœuds} est \textbf{égale à la somme des intensités} des courants électriques \textbf{qui en repartent}.

\vspace{3cm}
\begin{figure}[h!]
	\centering
	\begin{circuitikz}[scale=1.2, transform shape]
		\draw
			(0,-1) [battery1]
			to (0,3)
			to [R] (2,3)
			to [short,*-] (2,3)
			to [R] (2,1)
			to [R] (2,-1)
			to [short,*-] (2,-1)
			to [short] (0,-1)
			to [short] (0,0)
		;
		\draw
			(2,3)
			to [R] (4,3)
			to [short] (5,3)
			to [R] (5,-1)
			to [short] (5,-1)
			to [short] (2,-1)
			to [cground] (2,-2)
		;
	\end{circuitikz}
\end{figure}


\subsection{La loi des mailles}
Dans une \textbf{maille orientée}, la \textbf{somme des tensions} est \textbf{nulle}.\\
Pour additionner les tensions de la maille orientée, on affecte :
\begin{itemize}
	\item un signe "+" à une tension si sa flèche est dans le même sens que celui du parcours
	\item un signe "-" dans le sens contraire
\end{itemize}

Remarque : la tension aux bornes d'un fil est nulle.
%\vspace{3cm}

\section{Caractéristiques d'un dipôle}
La \textbf{caractéristique tension-courant} d'un dipôle est la représentation graphique de l'\textbf{évolution de la tension $U$ en fonction de l'intensité $I$} du courant électrique.\\
C'est la courbe d'équation $U=f(I)$\\
La caractéristique \textbf{courant-tension} d'un dipôle est la représentation graphique de l'\textbf{évolution de l'intensité $I$ du courant électrique en fonction de la tension $U$}.\\
C'est la courbe d'équation $I=g(U)$.

\subsection{Caractéristique d'un dipôle ohmique}
La caractéristique tension-courant $U=f(I)$ d'un dipôle ou du conducteur ohmique est une droite qui passe par l'origine, d'équation $U=R \times I$. La résistance $R$ est alors égale au coefficient directeur de la droite.\\
La tension $U$ et l'intensité $I$ du courant électrique sont donc liées par une relation de proportionnalité régie par le \textbf{Loi d'Ohm}.\\
D'après la \textbf{loi d'Ohm}, la tension $U$ aux bornes d'un conducteur ohmique est \textbf{égale} au \textbf{produit} de sa résistance $R$ et de l'intensité $I$ du courant électrique qui la traverse : 
\textbf{$U=R \times I$ avec $U$ en $V$, $R$ en $\Omega$ et $I$ en $A$}

%\vspace{3cm}

\subsection{Point de fonctionnement}
Lorsqu'un dipôle récepteur est branché aux bornes d'un générateur, un \textbf{courant de même intensité $I_f$} traverse les deux dipôle. La \textbf{tension $U_f$} à leurs bornes est également \textbf{la même}.\\
Ces deux conditions sont réalisées simultanément à l'intersection des caractéristiques des deux dipôles.\\
Le \textbf{point de fonctionnement d'un circuit, noté P($I_f$; $U_f$)}, est le \textbf{point d'intersection} des caractéristiques du générateur et du dipôle récepteur branché en série.\\

\section*{Fihe de révision}
\begin{tabular}{|m{.6\linewidth}|>{\centering}m{.15\linewidth}|>{\centering\arraybackslash}m{.15\linewidth}|}
	\hline
	\centering \textbf{Caractéristiques exigibles} & \textbf{Où la trouver} & \textbf{Maîtrise}\\
	\hline
	Exploiter la loi des mailles et la loi des n\oe uds dans un circuit électrique comportant au plus deux mailles. & & \\
	\hline
	Mesurer une tension et une intensité & &\\
	\hline
	Exploiter la caractéristique d'un dipôle électrique : point de fonctionnement, modélisation par la relation $U = f(t)$ ou $I = g(U)$. & &\\
	\hline
	Représenter et exploiter la caractéristique d'un dipôle et modéliser la caractéristique de ce dipôle à l'aide d'un langage de programmation. & & \\
	\hline
	Identifier une situation de proportionnalité & &\\
	\hline
	Citer des exemples de capteurs présents dans les objets de la vie quotidienne.& & \\
	\hline 
	Mesurer une grandeur physique à l'aide d'un capteur électrique résistif. Produire et utiliser une courbe d'étalonnage utilisant la résistance d'un système avec une grandeur d'intérêt (température, pression, intensité lumineuse, etc.) & & \\
	\hline
	Utiliser un dispositif avec micro-contrôleur et capteur. & &\\
	Exploiter une série de mesures, discuter de l'influence du protocole et/ou évaluer une incertitude-type pour comparer des résultats &&\\
	\hline
\end{tabular}
\end{document}